\chapter{Aufbau der Eingabedatei}
\label{kap:format-grundlagen}

Dieser Teil beschreibt das Eingabeformat vollständig. Er gilt gleichermaßen
für Dateien, die im Editor entstehen, und für Dateien, die man von Hand
schreibt oder aus einem Programm erzeugt.

\section{Die Datei als Ganzes}

Eine Eingabedatei ist ein JSON-Objekt. Anders als reines JSON darf sie
Kommentare enthalten: \code{//} bis zum Zeilenende und \code{/* ... */}
über mehrere Zeilen. Die Kodierung ist UTF-8. Ein Komma nach dem letzten
Element einer Liste oder eines Objekts ist, wie in JSON, ein Fehler.

\begin{center}
\begin{tabularx}{\textwidth}{@{}lLl@{}}
  \toprule
  Schlüssel & Inhalt & Abschnitt \\
  \midrule
  \code{title} & Titel, erscheint auf der Konsole und im Ergebnis & \\
  \code{parameters} & benannte Konstanten & \ref{sec:parameter} \\
  \code{functions} & eigene Funktionen für Ausdrücke & \ref{sec:funktionen} \\
  \code{tables} & Kennlinien aus Stützstellen & \ref{sec:tabellen} \\
  \code{models} & Modellparameter von Halbleitern & \ref{kap:format-modelle} \\
  \code{subcircuits} & Teilschaltungen & \ref{kap:format-teilschaltungen} \\
  \code{components} & die Bauteile der Schaltung & \ref{kap:format-bauteile} \\
  \code{variables} & Unbekannte eigener Gleichungen & \ref{kap:format-gleichungen} \\
  \code{equations} & eigene Gleichungen & \ref{kap:format-gleichungen} \\
  \code{events} & Ereignisse mit Zuweisungen & \ref{sec:ereignisse} \\
  \code{analyses} & die Analysen, in dieser Reihenfolge gerechnet & \ref{kap:format-analysen} \\
  \code{outputs} & auszugebende Größen & \ref{sec:ausgaenge} \\
  \code{options} & Toleranzen, Temperatur, Verfahren & \ref{sec:optionen} \\
  \code{nodeset} & Startwerte für den Arbeitspunkt & \ref{sec:optionen} \\
  \code{ic} & Anfangsspannungen für \code{"uic"} & \ref{sec:optionen} \\
  \code{schematic} & der Schaltplan des Editors (vom Simulator überlesen) & \ref{kap:editor-schaltplan} \\
  \bottomrule
\end{tabularx}
\end{center}

Eine Datei enthält \code{components}, \code{equations} oder beides.
Unbekannte Schlüssel – auf oberster Ebene, an Bauteilen, Modellen oder
Analysen – erzeugen eine Warnung und werden übergangen. So fallen
Tippfehler auf, ohne dass die Rechnung abbricht. Die Schlüssel von Bauteilen
und Modellen unterscheiden nicht zwischen Groß- und Kleinschreibung
(\code{"Value"} ist \code{"value"}).

\section{Zahlen}
\label{sec:zahlen}

Ein Zahlenwert kann sein:

\begin{itemize}
  \item eine JSON-Zahl: \code{1000}, \code{4.7e3}, \code{-0.5},
  \item eine Zeichenkette mit SI-Präfix und optionaler Einheit:
    \code{"4.7k"}, \code{"100nF"}, \code{"2.2M"}, \code{"10meg"},
  \item ein Ausdruck aus Zahlen und Parametern: \code{"2*R0"},
    \code{"1/(2*pi*R0*C0)"}.
\end{itemize}

\begin{center}
\begin{tabular}{@{}*{9}{c}@{}}
  \toprule
  \code{f} & \code{p} & \code{n} & \code{u}, \texttt{µ} & \code{m} &
  \code{k}, \code{K} & \code{M}, \code{meg} & \code{G} & \code{T} \\
  \midrule
  $10^{-15}$ & $10^{-12}$ & $10^{-9}$ & $10^{-6}$ & $10^{-3}$ &
  $10^{3}$ & $10^{6}$ & $10^{9}$ & $10^{12}$ \\
  \bottomrule
\end{tabular}
\end{center}

\begin{achtung}[Groß- und Kleinschreibung]
  Anders als in SPICE unterscheiden die Präfixe Groß- und Kleinschreibung:
  \code{"1M"} ist ein Megaohm, \code{"1m"} ein Milliohm. \code{meg} ist
  immer Mega. In SPICE-Text (Modellkarten, SPICE-Teilschaltungen) gilt
  dagegen die SPICE-Regel: \code{M} ist dort Milli.
\end{achtung}

Buchstaben nach dem Präfix gelten als Einheit und werden ignoriert:
\code{"100nF"} ist $10^{-7}$, \code{"1F"} ist 1 (ein Farad, nicht ein Femto).
In Ausdrücken ist nur das Präfix erlaubt (\code{"4.7k*2"}), keine Einheit.
Zahlen werden unabhängig von den Ländereinstellungen des Rechners gelesen;
das Dezimalzeichen ist immer der Punkt.

\section{Parameter}
\label{sec:parameter}

Parameter sind benannte Konstanten. Sie dürfen sich aufeinander beziehen,
die Reihenfolge spielt keine Rolle:

\begin{json}
"parameters": { "R0": "1k", "C0": "100n", "tau": "R0*C0", "fc": "1/(2*pi*tau)" }
\end{json}

Parameter können überall verwendet werden, wo eine Zahl steht – in
Bauteilwerten, Modellparametern, Analysen (\code{"tstop": "10*tau"}), in
Gleichungen und Ausgängen. Parameter-Sweeps, Monte Carlo,
Empfindlichkeiten und die Anpassung an Messdaten verändern Parameter; alles,
was von ihnen abhängt, wird für jeden Lauf neu berechnet. Wer eine Größe
später variieren will, schreibt sie deshalb am besten gleich als Parameter.

\section{Ausdrücke}
\label{sec:ausdruecke}

Ausdrücke stehen in Gleichungen, Verhaltensquellen, Ausgängen, Ereignissen,
Schaltern und überall, wo eine Zahl erwartet wird (dort nur mit
Parametern).

\begin{center}
\begin{tabularx}{\textwidth}{@{}P{3.3cm}L@{}}
  \toprule
  Element & Schreibweise \\
  \midrule
  Grundrechenarten & \code{+ - * /}, Potenz \code{^} oder \code{**}, Klammern \\
  Vergleiche & \code{< <= > >= == !=} (Ergebnis 1 oder 0) \\
  Logik & \code{&&}, \code{||}, \code{!} \\
  Fallunterscheidung & \code{if(bedingung, dann, sonst)} \\
  Funktionen & \code{sin cos tan asin acos atan atan2(y,x) sinh cosh tanh exp log ln log10 sqrt abs sign min(a,b) max(a,b) pow(a,b)} \\
  Konstanten & \code{pi} \\
  Zeit & \code{t} oder \code{time} \\
  Schaltungsgrößen & \code{v(knoten)}, \code{v(a,b)} $=$ \code{v(a)-v(b)}, \code{i(bauteil)} \\
  Variablen & Namen aus \code{variables}, Elemente \code{T[i]} von Vektoren \\
  Tabellen & \code{name(x)} für eine Kennlinie aus \code{tables} \\
  eigene Funktionen & \code{name(a, b)} aus \code{functions} \\
  Totzeit & \code{delay(x, td)} (Abschnitt~\ref{sec:totzeit}) \\
  Summe & \code{sum(i, von, bis, ausdruck)} (Abschnitt~\ref{sec:vektoren}) \\
  Zeitableitung & \code{der(ausdruck)}, nur in Gleichungen \\
  \bottomrule
\end{tabularx}
\end{center}

\code{log} und \code{ln} sind beide der natürliche Logarithmus; den
dekadischen liefert \code{log10}. Winkel sind im Bogenmaß.

Die Ableitungen für das Newton-Verfahren werden automatisch und exakt
gebildet. Unstetige Ausdrücke – Vergleiche, \code{if}, \code{abs},
\code{sign}, \code{min}, \code{max} – sind in Gleichungen und
Verhaltensquellen erlaubt: Die Stelle, an der ein Vergleich kippt, wird als
Ereignis exakt bestimmt (Abschnitt~\ref{sec:ereignisse}).

\section{Eigene Funktionen}
\label{sec:funktionen}

Wiederkehrende Ausdrücke lassen sich als Funktionen ablegen:

\begin{json}
"functions": {
  "clip(x, lim)":   "max(-lim, min(lim, x))",
  "core(i, L, is)": "L*is*tanh(i/is)"
}
\end{json}

Ein Aufruf wird durch den Rumpf ersetzt, mit den Argumenten anstelle der
Parameter. Funktionen dürfen Parameter der Datei und andere Funktionen
verwenden, sich aber nicht selbst aufrufen und keine eingebaute Funktion
überschreiben. Sie stehen in Gleichungen, Verhaltensquellen, Schaltern,
Ereignissen und Ausgängen zur Verfügung, nicht aber in reinen Zahlenwerten
(Bauteilwerte, \code{parameters}).

\section{Tabellen}
\label{sec:tabellen}

Kennlinien aus Stützstellen, etwa Messdaten oder Datenblattkurven, werden
unter \code{tables} definiert und in Ausdrücken wie eine Funktion mit einem
Argument aufgerufen:

\begin{json}
"tables": {
  "cell": { "x": [0, 0.3, 0.5, 0.6], "y": [3.0, 2.96, 2.38, 0], "interp": "cubic" },
  "ntc":  { "file": "ntc_curve.csv", "xcol": 0, "ycol": 1 }
},
"components": [
  { "type": "B", "name": "Bcell", "nodes": ["0", "p"], "i": "cell(v(p))" }
]
\end{json}

\begin{description}[style=nextline,leftmargin=2em]
  \item[Daten]
    \code{x} und \code{y} als Listen, oder \code{points} als Liste von
    Paaren \code{[x, y]}, oder \code{file}: eine Textdatei mit
    Zahlenspalten. Trenner sind Komma, Semikolon oder Leerzeichen; Zeilen,
    die nicht mit einer Zahl beginnen (Überschriften), werden übersprungen.
    Der Pfad gilt relativ zur Eingabedatei; \code{xcol} und \code{ycol}
    wählen die Spalten (gezählt ab 0). \code{x} muss streng steigen.
  \item[{\code{interp}}]
    \code{"linear"} (Standard) oder \code{"cubic"}. Die kubische
    Interpolation ist monoton (PCHIP): glatte Steigung, aber kein
    Überschwingen zwischen den Stützstellen. Für Kennlinien in Gleichungen
    und Quellen ist sie meist die bessere Wahl, weil das Newton-Verfahren
    dann eine stetige Ableitung sieht.
  \item[{\code{extrap}}]
    Außerhalb der Daten wird der Randwert gehalten (\code{"hold"},
    Standard) oder mit der Randsteigung fortgesetzt (\code{"linear"}).
\end{description}

\subsection*{Kennfelder über zwei Größen}

Ein Kennfeld, etwa die Flussverkettung über Strom und Position aus einer
Feldrechnung, wird mit \code{z} oder \code{zcol} angelegt und mit zwei
Argumenten aufgerufen, \code{name(x, y)}:

\begin{json}
"tables": {
  "psi": { "x": [0, 1, 2], "y": [0, 5e-3, 1e-2],
           "z": [[0, 0, 0], [0.01, 0.006, 0.004], [0.018, 0.011, 0.007]],
           "interp": "cubic" },
  "map": { "file": "kennfeld.csv", "xcol": 0, "ycol": 1, "zcol": 2 }
},
"components": [
  { "type": "B", "name": "L1", "nodes": ["a", "0"], "phi": "psi(i(L1), z)" }
]
\end{json}

\code{z[i][j]} gehört zu \code{(x[i], y[j])}. Eine Datei enthält einen Punkt
je Zeile in beliebiger Reihenfolge und muss das ganze Gitter abdecken.
\code{x} und \code{y} müssen streng steigen. \code{"linear"} ist bilinear;
\code{"cubic"} ist ein Tensorprodukt kubischer Hermite-Interpolanten mit
stetiger Steigung und exakten Ableitungen für das Newton-Verfahren, aber
(anders als bei der Kennlinie) nicht monoton. \code{extrap} gilt für beide
Achsen.

\section{Knoten}

Knotennamen sind beliebige Zeichenketten oder ganze Zahlen. Masse heißt
\code{"0"} oder \code{"gnd"}. Jede Schaltung braucht einen Massebezug, und
jeder Knoten braucht einen Gleichstrompfad nach Masse
(Abschnitt~\ref{sec:schwebend}). In Ausdrücken (\code{v(...)}) sind nur
Knotennamen ohne Komma und Klammern verwendbar; einfache Namen wie
\code{out}, \code{vdd}, \code{n12} sind deshalb am praktischsten.

Größen innerhalb von Teilschaltungen heißen von außen
\code{<Instanz>.<Name>}, etwa \code{v(X1.mid)} oder \code{i(U1.V1)}
(Kapitel~\ref{kap:format-teilschaltungen}).
